\subsection{Congruências modulares}

Sejam $a, b$ e $n$ inteiros dados. Queremos achar os $x$ inteiros em $[0, n-1]$ tais que
$$ a \cdot x \equiv b ( \mod n)$$
Seja $g = \gcd(a, n)$. Seja $a=a'g$, $n=n'g$. Se $g$ não dividir $b$, não existem soluções. Seja $b=gb'$. A equação original implica
$$a' \cdot x \equiv b' (\mod n')$$
Como $\gcd(a', n') = 1$, então só existe uma solução. Então, ao todo, existem $g$ soluções dadas por
$$x_i \equiv x' + i \cdot n' ( \mod n), \, i = 0 \dots g-1$$

Também é possível resolvê-la utilizando o algoritmo estendido de euclides com a equação
$$a \cdot x + n \cdot k = b$$

em que $x$ e $k$ são desconhecidos.

\subsubsection{Teorema Chinês dos Restos}

Sejam $m_1, \dots, m_k$ coprimos. Sejam $a_i$ inteiros. Queremos achar $a$ tal que $a \equiv a_i (\mod m_i)$ para todo $1 \leq i \leq k$. 

A solução é única mod $m_1 m_2 \dots m_k$.

Para construi-la, seja $M_i = \prod_{i \neq j} m_j$ e $N_i = M_i^{-1} \mod m_i$. Então, módulo $m_1 \dots m_k$, a solução é $$a \equiv \sum\limits_{i=1}^k a_i M_i N_i$$

\begin{lstlisting}[language=c++, breaklines=true]
struct Congruence {
    long long a, m;
};

long long chinese_remainder_theorem(vector<Congruence> const& congruences) {
    long long M = 1;
    for (auto const& congruence : congruences) {
        M *= congruence.m;
    }

    long long solution = 0;
    for (auto const& congruence : congruences) {
        long long a_i = congruence.a;
        long long M_i = M / congruence.m;
        long long N_i = mod_inv(M_i, congruence.m);
        solution = (solution + a_i * M_i % M * N_i) % M;
    }
    return solution;
}
\end{lstlisting}

\subsubsection{Algoritmo de Garner}

É possível usar o teorema chinês dos restos para representar um número muito gigante $a$ utilizando apenas os $a_i$. Utilizamos a representação
$$a = x_1 + x_2 p_1 + x_3 p_1 p_2 + \dots + x_k p_1 \dots p_{k-1}$$
com $x_i \in [0, p_i)$

O que o algoritmo de Garner faz é computar os dígitos $x_i$. É possível utilizar esse algoritmo para calcular em $O(k^2)$ os $x_i$. Eis uma implementação \textbf{em java}:

\begin{lstlisting}[language=java, breaklines=true]
final int SZ = 100;
int pr[] = new int[SZ];
int r[][] = new int[SZ][SZ];

void init() {
    for (int x = 1000 * 1000 * 1000, i = 0; i < SZ; ++x)
        if (BigInteger.valueOf(x).isProbablePrime(100))
            pr[i++] = x;

    for (int i = 0; i < SZ; ++i)
        for (int j = i + 1; j < SZ; ++j)
            r[i][j] =
                BigInteger.valueOf(pr[i]).modInverse(BigInteger.valueOf(pr[j])).intValue();
}

class Number {
    int a[] = new int[SZ];

    public Number() {
    }

    public Number(int n) {
        for (int i = 0; i < SZ; ++i)
            a[i] = n % pr[i];
    }

    public Number(BigInteger n) {
        for (int i = 0; i < SZ; ++i)
            a[i] = n.mod(BigInteger.valueOf(pr[i])).intValue();
    }

    public Number add(Number n) {
        Number result = new Number();
        for (int i = 0; i < SZ; ++i)
            result.a[i] = (a[i] + n.a[i]) % pr[i];
        return result;
    }

    public Number subtract(Number n) {
        Number result = new Number();
        for (int i = 0; i < SZ; ++i)
            result.a[i] = (a[i] - n.a[i] + pr[i]) % pr[i];
        return result;
    }

    public Number multiply(Number n) {
        Number result = new Number();
        for (int i = 0; i < SZ; ++i)
            result.a[i] = (int)((a[i] * 1l * n.a[i]) % pr[i]);
        return result;
    }

    public BigInteger bigIntegerValue(boolean can_be_negative) {
        BigInteger result = BigInteger.ZERO, mult = BigInteger.ONE;
        int x[] = new int[SZ];
        for (int i = 0; i < SZ; ++i) {
            x[i] = a[i];
            for (int j = 0; j < i; ++j) {
                long cur = (x[i] - x[j]) * 1l * r[j][i];
                x[i] = (int)((cur % pr[i] + pr[i]) % pr[i]);
            }
            result = result.add(mult.multiply(BigInteger.valueOf(x[i])));
            mult = mult.multiply(BigInteger.valueOf(pr[i]));
        }

        if (can_be_negative)
            if (result.compareTo(mult.shiftRight(1)) >= 0)
                result = result.subtract(mult);

        return result;
    }
}
\end{lstlisting}

\subsubsection{Fatorial módulo $p$}

Dá pra precalcular em $O(p)$ os fatoriais e em $O(\log_p n)$ o fatorial de um determinado número. Isso usa $O(p)$ de memória.

\begin{lstlisting}[language=c++, breaklines=true]
int factmod(int n, int p) {
    vector<int> f(p);
    f[0] = 1;
    for (int i = 1; i < p; i++)
        f[i] = f[i-1] * i % p;

    int res = 1;
    while (n > 1) {
        if ((n/p) % 2)
            res = p - res;
        res = res * f[n%p] % p;
        n /= p;
    }
    return res;
}
\end{lstlisting}

Truque de mestre! Se $p \approx 10^9$, dá pra hardcodar um vetor fatmod tal que fatmod[i] é $(Si)!$ módulo $p$ em que $S$ é um tamanho de bloco. Dá pra usar $S=10^6$ pro vetor ter $10^3$ caras e isso ser algo colocável no código.

Então ficamos com:

\begin{lstlisting}[language=c++, breaklines=true]
const long long int mod = 1e9+7;
const long long int maxn = 1e9;
const int tempo = 1e6;
long long int fatmod[mod/tempo+1] = { /* com outro programa calcular os numeros e copiar e colar aqui*/};
 
 
long long int computa_fatmod(int n) {
    if (n >= mod) return 0;
    long long int atual = fatmod[n/tempo];
    for (int i = (n/tempo)*tempo+1; i <= n; i++) {
        atual *= i;
        atual %= mod;
    }
    return atual;
}
\end{lstlisting}

\subsubsection{Fórmula de Legendre}

Qual o expoente de $p$ na fatoração prima de $n!$? Sem calcular $n!$, podemos calcular isso em $O( \log_p n)$ com a seguinte fórmula:
$$ v_p(n!) = \sum\limits_{i=1}^\infty \left\lfloor \dfrac{n}{p^i} \right\rfloor$$

\begin{lstlisting}[language=c++, breaklines=true]
int multiplicity_factorial(int n, int p) {
    int count = 0;
    do {
        n /= p;
        count += n;
    } while (n);
    return count;
}
\end{lstlisting}

\subsubsection{Logaritmo Discreto}

Inteiro $x$ satisfazendo 
$$a^x \equiv b ( \mod m)$$
para inteiros $a, b, m$ dados. Será resolvido com o algoritmo \textbf{baby-step giant-step} de complexidade $O(\sqrt{m})$. Assumiremos que $0^0 \equiv 1$.

Forma mais simples sem usar unordered map roda em $O(\sqrt{m} \log m)$ (assumiremos nas duas versões $\gcd(a, m)=1$):

\begin{lstlisting}[language=c++, breaklines=true]
int powmod(int a, int b, int m) {
    int res = 1;
    while (b > 0) {
        if (b & 1) {
            res = (res * 1ll * a) % m;
        }
        a = (a * 1ll * a) % m;
        b >>= 1;
    }
    return res;
}

int solve(int a, int b, int m) {
    a %= m, b %= m;
    int n = sqrt(m) + 1;
    map<int, int> vals;
    for (int p = 1; p <= n; ++p)
        vals[powmod(a, p * n, m)] = p;
    for (int q = 0; q <= n; ++q) {
        int cur = (powmod(a, q, m) * 1ll * b) % m;
        if (vals.count(cur)) {
            int ans = vals[cur] * n - q;
            return ans;
        }
    }
    return -1;
}
\end{lstlisting}

Usando unordered map:

\begin{lstlisting}[language=c++, breaklines=true]
// Returns minimum x for which a ^ x % m = b % m, a and m are coprime.
int solve(int a, int b, int m) {
    a %= m, b %= m;
    int n = sqrt(m) + 1;

    int an = 1;
    for (int i = 0; i < n; ++i)
        an = (an * 1ll * a) % m;

    unordered_map<int, int> vals;
    for (int q = 0, cur = b; q <= n; ++q) {
        vals[cur] = q;
        cur = (cur * 1ll * a) % m;
    }

    for (int p = 1, cur = 1; p <= n; ++p) {
        cur = (cur * 1ll * an) % m;
        if (vals.count(cur)) {
            int ans = n * p - vals[cur];
            return ans;
        }
    }
    return -1;
}
\end{lstlisting}

Caso em que $\gcd(a, m) \neq 1$. Se $g$ não divide $b$, então não há solução. Se $g | b$, então tomando $a = g \alpha$, $b = g \beta$, $m = g \nu$. Daí, precisamos resolver $\alpha a^{x-1} \equiv \beta \mod \nu$. Mas o nosso algoritmo pode ser facilmente estendido para $k a^x \equiv b$:

\begin{lstlisting}[language=c++, breaklines=true]
// Returns minimum x for which a ^ x % m = b % m.
int solve(int a, int b, int m) {
    a %= m, b %= m;
    int k = 1, add = 0, g;
    while ((g = gcd(a, m)) > 1) {
        if (b == k)
            return add;
        if (b % g)
            return -1;
        b /= g, m /= g, ++add;
        k = (k * 1ll * a / g) % m;
    }

    int n = sqrt(m) + 1;
    int an = 1;
    for (int i = 0; i < n; ++i)
        an = (an * 1ll * a) % m;

    unordered_map<int, int> vals;
    for (int q = 0, cur = b; q <= n; ++q) {
        vals[cur] = q;
        cur = (cur * 1ll * a) % m;
    }

    for (int p = 1, cur = k; p <= n; ++p) {
        cur = (cur * 1ll * an) % m;
        if (vals.count(cur)) {
            int ans = n * p - vals[cur] + add;
            return ans;
        }
    }
    return -1;
}
\end{lstlisting}

\subsubsection{Raízes primitivas}

$g$ é dita uma raiz primitiva módulo $n$ se todo número coprimo a $n$ é congruente a uma potência de $g$ módulo $n$. Se $g$ é uma raiz primitiva e $g^k \equiv a$, então $k$ é dito o índice ou logaritmo discreto de $a$ na base $g$.

Uma raiz primitiva módulo $n$ existe se e só se:

\begin{enumerate}
    \item $n=1,2,4$ ou
    \item $n = p^k$, $p$ primo ímpar, $k \geq 1$
    \item $n = 2p^k$, $p$ primo ímpar, $k \geq 1$
\end{enumerate}

Se existe alguma raiz primitiva módulo $n$, então existem $\phi(\phi(n))$ delas. O seguinte código acha um gerador em $O(Ans \cdot \log \phi(n) \cdot \log n)$. Sabe-se que $g \in O(\log^6 p)$. Há um detalhe: o código abaixo \textbf{só funciona} para $p$ primo. Mas pra funcionar pra qualquer número basta trocar a definição da variável phi:

\begin{lstlisting}[language=c++, breaklines=true]
int powmod (int a, int b, int p) {
    int res = 1;
    while (b)
        if (b & 1)
            res = int (res * 1ll * a % p),  --b;
        else
            a = int (a * 1ll * a % p),  b >>= 1;
    return res;
}

int generator (int p) {
    vector<int> fact;
    int phi = p-1,  n = phi;
    for (int i=2; i*i<=n; ++i)
        if (n % i == 0) {
            fact.push_back (i);
            while (n % i == 0)
                n /= i;
        }
    if (n > 1)
        fact.push_back (n);

    for (int res=2; res<=p; ++res) {
        bool ok = true;
        for (size_t i=0; i<fact.size() && ok; ++i)
            ok &= powmod (res, phi / fact[i], p) != 1;
        if (ok)  return res;
    }
    return -1;
}
\end{lstlisting}

\subsubsection{Raiz discreta}

Dado um primo $n$ e dois inteiros $a, k$, achar todos os $x$ tais que $x^k \equiv a \mod n$.

\textbf{Solução:} achar uma raiz primitiva $g$ módulo $n$ (em $O(Ans \cdot \log^2 n)$). Depois, se $x \equiv g^y$, basta resolver $(g^k)^y \equiv a$ em $y$ usando o logaritmo discreto.

\prob{1}{congmod1} Encontrar todas as soluções a partir de uma.

\textbf{Solução:} suponha que conhecemos a solução $x_0 = g^{y_0} \mod n$. Então, todas as soluções são da forma
$$x \equiv g^{y_0 + i \dfrac{\phi(n)}{\gcd(k, \phi(n))}}$$
com $i$ variando em $\mathbb{Z}$.

O seguinte código mostra todas as soluções:

\begin{lstlisting}[language=c++, breaklines=true]
int gcd(int a, int b) {
    return a ? gcd(b % a, a) : b;
}

int powmod(int a, int b, int p) {
    int res = 1;
    while (b > 0) {
        if (b & 1) {
            res = res * a % p;
        }
        a = a * a % p;
        b >>= 1;
    }
    return res;
}

// Finds the primitive root modulo p
int generator(int p) {
    vector<int> fact;
    int phi = p-1, n = phi;
    for (int i = 2; i * i <= n; ++i) {
        if (n % i == 0) {
            fact.push_back(i);
            while (n % i == 0)
                n /= i;
        }
    }
    if (n > 1)
        fact.push_back(n);

    for (int res = 2; res <= p; ++res) {
        bool ok = true;
        for (int factor : fact) {
            if (powmod(res, phi / factor, p) == 1) {
                ok = false;
                break;
            }
        }
        if (ok) return res;
    }
    return -1;
}

// This program finds all numbers x such that x^k = a (mod n)
int main() {
    int n, k, a;
    scanf("%d %d %d", &n, &k, &a);
    if (a == 0) {
        puts("1\n0");
        return 0;
    }

    int g = generator(n);

    // Baby-step giant-step discrete logarithm algorithm
    int sq = (int) sqrt (n + .0) + 1;
    vector<pair<int, int>> dec(sq);
    for (int i = 1; i <= sq; ++i)
        dec[i-1] = {powmod(g, i * sq * k % (n - 1), n), i};
    sort(dec.begin(), dec.end());
    int any_ans = -1;
    for (int i = 0; i < sq; ++i) {
        int my = powmod(g, i * k % (n - 1), n) * a % n;
        auto it = lower_bound(dec.begin(), dec.end(), make_pair(my, 0));
        if (it != dec.end() && it->first == my) {
            any_ans = it->second * sq - i;
            break;
        }
    }
    if (any_ans == -1) {
        puts("0");
        return 0;
    }

    // Print all possible answers
    int delta = (n-1) / gcd(k, n-1);
    vector<int> ans;
    for (int cur = any_ans % delta; cur < n-1; cur += delta)
        ans.push_back(powmod(g, cur, n));
    sort(ans.begin(), ans.end());
    printf("%d\n", ans.size());
    for (int answer : ans)
        printf("%d ", answer);
}
\end{lstlisting}